\chapter{Discussion}
\label{ch:discussion}

\section{What the results mean}
\label{sec:d-meaning}

The five results of Chapter~\ref{ch:eval} compose into a single claim with a sharp
boundary. On the repeatable head of on-device tasks, a 2B reflexer with a
distilled pattern matches a $17\times$ larger free-form thinker, because the
reflexer does not solve the task---it replays a solution found once. The boundary
is the distinction the thesis makes precise: model scale buys \emph{discovery},
the search for which namespace and key a request maps to, but it does not buy
\emph{convention}, the device-specific meaning of a value. A 35B thinker explores
the API and finds the right key, then guesses the encoding and is wrong about half
the time, making the \emph{same} mistake every time it is asked
(Section~\ref{sec:p4}). A distilled pattern carries that convention, so a small
model applying it succeeds where the large model reasoning from scratch fails.

This reframes a question usually posed as ``how capable a model do we need'' into
``where does the required knowledge live.'' For discovery, the knowledge is in the
API and a capable model can find it. For convention, the knowledge is in neither
the API nor the model's parameters at any scale tested; it exists only in a
distilled artifact. Once that is seen, the economics follow directly
(Section~\ref{sec:economics}): pay a strong teacher once to externalize the
convention, and a cheap student applies it forever. The teacher-quality result
sharpens the recipe---because the student faithfully executes whatever convention
the teacher encoded, a weak teacher bakes in wrong priors, so the one-time
distillation step should use the strongest available teacher
(Section~\ref{sec:p5}).

\section{Implications for productization}
\label{sec:d-product}

The architecture suggests a clear product shape and an equally clear set of
blockers. The defensible first version is a \emph{deterministic-intent
accelerator}: the high-frequency, unambiguous head---media and device control,
smart-home toggles, timers, system actions---served on-device by the reflexer at
near-100\% binding, with the thinker as a safety net for everything else. This is
the slice where the value is highest and the risk is lowest, and it coincides with
the on-device sweet spot that the Pixel and Apple use-case mapping identifies
(Section~\ref{sec:coverage}).

The blockers are ranked by how much they gate a shippable system. \emph{Routing
reliability} is first: end-to-end accuracy is the product of routing and binding,
binding is solved, and routing at 66--85\% is below what users tolerate for an
assistant (Section~\ref{sec:routing}); a misroute to the wrong action is the
dangerous failure, worse than a refusal. \emph{Cold-start data and authoring cost}
is second: each skill needs a pattern set, and patterns must be distilled and,
ideally, kept current as device APIs change; the distillation pipeline is what
makes this scale. \emph{Composite slot extraction} is third: multi-step routines
that need values produced at run time are the high-value automation surface and
remain the reflexer's weak point. None of these is a flaw in the executioner idea;
they are the engineering distance between a validated mechanism and a product.

\section{Threats to validity}
\label{sec:d-threats}

Three caveats must travel with the headline numbers, and the thesis states them
rather than burying them.

First, the comparison is \emph{conditional on a constrained task}. The reflexer is
handed the command template and only fills slots, which is strictly easier than
free-form generation. The claim is therefore not that a 2B is more capable than a
35B; it is that on the distillable slice a constrained 2B executor reaches the
same binding accuracy at a fraction of the cost. Read as a capability claim it
would be wrong; read as a systems claim it is what the numbers support.

Second, \emph{coverage is limited to distilled intents}. The 100\% binding figure
is a ceiling that applies only when a pattern exists and routing has selected it.
For intents with no pattern, the system falls back to the thinker, and the
reflexer contributes nothing. The approach is a head strategy by construction; its
value depends on the workload actually having a repeatable head, which the
coverage mapping argues it does.

Third, the strongest results rest on a \emph{small number of comparisons on
purpose-built benchmarks} with specific models. SET-2 is realistic but is a model
of the Android settings surface, not the full device; the 35B is one large model,
not a survey of large models; and two runs of the cold-start reflexer disagreed by
one task (19 versus 20 of 25), which the thesis reports rather than reconciles. The
qualitative pattern---scale fixes discovery but not convention, distillation fixes
both, teacher quality propagates---is robust across the experiments, but the exact
percentages should be read as evidence for that pattern rather than as universal
constants.

\section{Limitations}
\label{sec:d-limits}

Beyond the threats to validity, four limitations bound the work. The router is
evaluated but not solved; it is named as the open bottleneck rather than fixed.
Composite slot-filling is identified as weak but not repaired. Both teachers,
cloud and local, miss the genuinely undiscoverable conventions (the version- or
OEM-specific values), which no cold-start teacher can recover and which require
observing the real device. And the reflexer's residual application errors---using
an enum label instead of its code, or mis-doing brightness arithmetic---point to a
value-transform step that is currently left to the model. Each of these limitations
maps onto a concrete next step rather than an open research question, which the
next chapter sets out.
